\section{Optimizing and Executing a Mixture-of-Models}
\label{sec:optimization}

MoM extends the optimization surface beyond constituent weights to the
distribution of computation around independently versioned models. A lifecycle
engine can optimize that distribution while preserving one MoM identity across
training, evaluation, and serving. Accordingly, MoM optimization is a
full-stack co-design problem. Changing the trace policy changes the demand
presented to the serving fleet; changing placement, batching, or pool capacity
changes the feasible trace distribution. Logical and physical coordinates must
therefore be optimized jointly while remaining separately versioned for
portability and attribution \citep{chen_2026_wrp}. Figure~\ref{fig:system-codesign}
makes this coupling explicit.

\begin{figure}[H]
  \centering
  \resizebox{\linewidth}{!}{\begin{tikzpicture}[
  x=1cm,y=1cm,font=\sffamily,>=Latex,
  flow/.style={-{Latex[length=1.9mm]},draw=MoMGray,line width=0.68pt},
  feedback/.style={-{Latex[length=1.8mm]},draw=MoMGray,dashed,line width=0.62pt},
  panel/.style={draw=MoMBorder,fill=white,rounded corners=3pt,
    line width=0.62pt,minimum width=3.45cm,minimum height=1.42cm,
    inner xsep=5pt,inner ysep=4pt,align=center,text=MoMText},
  logical/.style={panel,draw=MoMDark,fill=MoMBlueLight!32,line width=0.78pt},
  physical/.style={panel,draw=MoMDark,fill=MoMBlue!72,line width=0.78pt},
  outcome/.style={draw=MoMBorder,fill=MoMDark,text=white,rounded corners=2.6pt,
    line width=0pt,minimum width=7.55cm,minimum height=0.82cm,
    inner xsep=6pt,inner ysep=3pt,align=center},
  meta/.style={font=\sffamily\fontsize{6.5}{7.2}\selectfont,text=MoMMuted}
]

\newcommand{\codesigntitle}{\bfseries\fontsize{8.2}{9.0}\selectfont}
\newcommand{\codesignmeta}{\fontsize{6.5}{7.2}\selectfont\color{MoMMuted}}
\newcommand{\codesignlightmeta}{\fontsize{6.5}{7.2}\selectfont\color{white}}

\node[panel] (workload) at (2.15,2.56)
  {{\codesigntitle Workload envelope}\\[-1pt]
   {\codesignmeta requests \textbullet\ preferences \textbullet\ constraints}};

\node[logical] (logical) at (7.98,2.56)
  {{\codesigntitle Logical MoM}\\[-1pt]
   {\codesignmeta signals \textbullet\ policy \textbullet\ typed traces}};

\node[physical] (physical) at (13.81,2.56)
  {{\codesigntitle Physical realization}\\[-1pt]
   {\codesignmeta binding \textbullet\ model pool \textbullet\ fleet}};

\draw[flow] (workload) -- node[meta,above] {shapes demand} (logical);
\draw[flow] (logical) -- node[meta,above] {requests capacity} (physical);
\draw[flow] (physical.south west) to[bend left=16]
  node[meta,below] {constrains feasible traces} (logical.south east);

\node[outcome] (outcome) at (7.98,0.68)
  {{\codesigntitle Realized trace and outcomes}\\[-1pt]
   {\codesignlightmeta quality \textbullet\ validity \textbullet\ latency
   \textbullet\ cost \textbullet\ failures}};

\draw[flow] (logical.south) -- (outcome.north);
\draw[flow] (physical.south) |- (outcome.east);
\draw[feedback] (outcome.west) -| node[meta,pos=0.33,below] {observe} (workload.south);

\end{tikzpicture}
}
  \caption{\textbf{MoM as system-level intelligence.} Workload, preferences,
  and constraints shape the useful trace distribution; the logical MoM shapes
  demand on constituent models; and bindings, placement, and fleet state bound
  feasible execution. Realized quality, validity, latency, cost, capacity, and
  failures feed back into both logical optimization and physical planning.}
  \label{fig:system-codesign}
\end{figure}

\subsection{Training the trace distribution}

Let $\mathcal{M}_{\theta_{\mathrm{log}}}$ denote a parameterized logical
specification. We write the lifecycle optimization coordinates as
\begin{equation}
 \Theta=(\theta_{\mathrm{log}},\beta),\qquad
 \theta_{\mathrm{log}}=(\phi,\omega,\psi,\rho,\eta),\qquad
 \mathcal{B}=\mathcal{B}_{\beta},\quad
 L=\operatorname{Resolve}(\mathcal{M}_{\theta_{\mathrm{log}}},
                           \mathcal{B}_{\beta},e).
 \label{eq:trainable-surface}
\end{equation}
$\phi$ parameterizes learned signals, $\omega$ their projections, and $\psi$
the action controller. The discrete coordinate $\rho$ selects catalog members,
topology, prompts, thresholds, contracts, and resource bounds already represented
in $\mathcal{M}$; it does not introduce a second semantic authority. Optional
$\eta$ denotes adapters or constituent weights that the optimizer may update.
The deployment coordinate $\beta$ controls binding, placement, replication, and
batching. Promoting $\theta_{\mathrm{log}}$ creates a new logical model version;
changing $\beta$ creates a new binding, after which resolution produces a new
lock.

Different coordinates call for different optimizers. Signals and routes can be
learned from supervision; rankings and escalation from preferences; thresholds
from logged outcomes or bandit feedback; pools and topologies through discrete
search; and prompts and budgets through black-box optimization. The resulting
policy may be distilled into a small controller. Any adapter or
constituent-model fine-tuning must be reported separately.

This definition does not turn every configuration edit into training. A MoM
training procedure must identify its experience source, optimization variables,
held-out objective, constraints, and produced artifact or binding version.
Constituent weights may remain frozen, be tuned independently, or be co-trained
when access and interfaces permit. Online adaptation must additionally bound
exploration and preserve policy invariants.

For a training set, the engine minimizes task loss and soft resource outcomes
subject to the hard admissibility conditions in Equation~\ref{eq:admissible}:
\begin{equation}
\begin{aligned}
 \min_{\Theta}\quad &\frac{1}{N}\sum_{i=1}^{N}
 \mathbb{E}_{\tau\sim p_{\mathcal{R}_{\Theta}}
 (\cdot\mid x_i,h_i,u_i,e_i)}
 \left[\ell(\bar y_i,y_i^*)+\lambda(u_i)^{\top}c(\tau,e_i)\right]\\
 \text{s.t.}\quad &
 \operatorname{supp}p_{\mathcal{R}_{\Theta}}
 (\cdot\mid x_i,h_i,u_i,e_i)\subseteq\mathcal{T}_{\mathrm{adm}}
 \quad\forall i.
\end{aligned}
 \label{eq:mom-training}
\end{equation}
Any allowed statistical violation rate must be stated explicitly and estimated
on held-out data; a soft penalty does not enforce an exact privacy or residency
rule.

Training produces a candidate version or asset, which evaluation compares with
the incumbent under a fixed protocol. Promotion is explicit, and serving never
silently rewrites the active specification. Bounded online adaptation remains
trace-visible and reversible.

\subsection{Capacity is heterogeneous and budgeted}
\label{sec:capacity}

Parameter count is a poor capacity measure: it may be unknown, inactive models
consume no generation compute, and similarly sized models can share errors.
We distinguish \emph{catalog capacity} (capability coverage), \emph{decision
capacity} (complementarity exploitable under a budget), and \emph{serving
capacity} (throughput, memory, concurrency, and failure envelope).

For pure selection and a per-example utility $q_i(x)$ for model $v_i$---with
$h$, $u$, and $e$ suppressed for notation---two diagnostic ceilings are
\begin{equation}
 Q_{\mathrm{single}}=\max_i\mathbb{E}_{x}[q_i(x)],\qquad
 Q_{\mathrm{oracle}}=\mathbb{E}_{x}\left[\max_i q_i(x)\right].
 \label{eq:oracle}
\end{equation}
Their gap $\Delta_{\mathrm{comp}}=Q_{\mathrm{oracle}}-Q_{\mathrm{single}}$
measures selection headroom. Let $Q_{\mathcal{M}}$ be the expected utility of a
pure-selection MoM on the same population and under the same budget. For a
positive gap, controller realization is
\begin{equation}
 \eta_{\mathrm{ctrl}}=
 \frac{Q_{\mathcal{M}}-Q_{\mathrm{single}}}
      {Q_{\mathrm{oracle}}-Q_{\mathrm{single}}}.
 \label{eq:controller-realization}
\end{equation}
Fusion may exceed this oracle because it combines multiple responses; report
that excess separately as composition gain. A large catalog provides little
decision capacity when constituent errors are strongly correlated.

Capacity scaling requires a predeclared matched active-work envelope over
calls, tokens, cost, or latency, with all resource outcomes reported.
Increasing catalog size under that envelope tests conditional capacity;
increasing catalog and active work only tests a larger system.

Let $\mathcal{F}(\mathcal{V},\bar b)$ be the feasible MoM policies over catalog
$\mathcal{V}$ under fixed budget $\bar b$ and fixed constraints, and let
\begin{equation}
 U^*(\mathcal{V},\bar b)=
 \sup_{\pi'\in\mathcal{F}(\mathcal{V},\bar b)}
 \mathbb{E}[R-\lambda^{\top}c].
 \label{eq:catalog-monotonicity}
\end{equation}
If $\mathcal{V}\subseteq\mathcal{V}'$ and every policy feasible over
$\mathcal{V}$ remains feasible over $\mathcal{V}'$, then
$U^*(\mathcal{V}',\bar b)\geq U^*(\mathcal{V},\bar b)$ because a policy over the larger
catalog can assign zero probability to every added model. This monotonicity
belongs to oracle capacity, not to a learned controller. Realized utility may
fall as a catalog grows when controller regret, exploration, or coordination
overhead increases faster than the new complementarity. Separating the benefit
of added complementarity from the cost of controller regret is the central
scaling question.

Let the expected activation of logical model $i$ per request be
\begin{equation}
 \Lambda_i=\mathbb{E}_{\tau}\left[
   \sum_{j\in N}\mathbf{1}[a_j=\textsc{Call}(v_i,\cdot)]\right].
 \label{eq:model-load}
\end{equation}
At request rate $\nu$, $\nu\Lambda_i$ determines its mean call rate. Models differ
in quality, price, limits, and capacity, so uniform $\Lambda_i$ is not a desirable
default.  Balancing should target the highest feasible utility under capacity
and tail-latency constraints, not equal traffic.

\subsection{Collapse and failure modes}

MoM couples statistical and systems failures. \textbf{Pool collapse} is harmful
when concentrated traffic excludes a better feasible policy; route entropy
alone cannot distinguish such collapse from correctly rejecting dominated
models. Diagnose it against $Q_{\mathrm{single}}$, oracle headroom, and realized
utility. \textbf{False complementarity} arises from leakage,
judge bias, or small samples and requires held-out, per-item error correlation.
\textbf{Policy dead zones} leave reachable signal states without valid traces;
static coverage and fail-closed tests expose them.

\textbf{Hidden escalation} underreports cost by omitting retries and fallbacks.
\textbf{Handoff loss} discards context through truncation, poor prompt
translation, or lossy summaries. \textbf{Contract collapse} masks invalid
outputs through unreported repair. Evaluations should therefore trace every
attempt and report first-pass validity, repairs, and terminal failures.

Finally, \textbf{provider drift and capability aliasing} break assumed
substitutability; locks, probes, and continuous evaluation detect and constrain
their effects. \textbf{Placement skew} overloads a pool or violates locality.
Model-generated plans introduce \textbf{unbounded control} unless the engine restricts them to declared
references and enforces $\Gamma$.

\subsection{Compile once, schedule under current state}

Before serving, the proposed engine type-checks the graph, validates contracts
and bounds, and resolves requirements. Per request it evaluates signals, filters
inadmissible actions, applies the preference-conditioned policy, and schedules
the topology.  Parallel latency follows the critical path.  Timeouts, partial
joins, fallbacks, caching, batching, and retries are declared semantics.

The engine emits, alongside each response, a trace linked to the MoM artifact
and resolution lock for attribution, replay, accounting, optimization, and
incident analysis. When policy permits, privacy-preserving redaction should
retain the minimum structural metadata needed to reconstruct the decision.

\paragraph{Reference-engine scope.}
vLLM Semantic Router grounds this execution contract without assuming access
to constituent weights or a jointly differentiable trainer.
Section~\ref{sec:reference-engine} audits the implemented subset and the
lifecycle boundary proposed here.
